\chapter{Dynamic Graph Algorithms}
\label{ch:dynamic}

In the preceding chapters, every algorithm operated on a fixed static graph: whether finding
shortest paths, max-flow, or approximation guarantees, the graph held still while computation ran.
But real graphs rarely hold still. A social network gains and loses friendships; a road map closes a
bridge for repair; a circuit simulator toggles wires; a chess-like search adds and removes pieces on
a board encoded as a graph. Recomputing from scratch after each change is the obvious thing to do,
and for a single static query it is exactly right. The moment changes come in a long stream,
recomputation is a disaster: an $\Oh(m)$ rebuild after each of $q$ updates costs $\Oh(qm)$, and for
$q\approx m$ that is quadratic.

\emph{Dynamic graph algorithms} ask the sharper question. Given a graph that
changes by small local edits -- insert an edge, delete an edge, sometimes change
a weight -- maintain some property (connectivity, a spanning forest, the MST
weight, the number of components, reachability, shortest-path distances) so that
each \textbf{update} and each \textbf{query} is fast, ideally polylogarithmic.
The art is to store just enough auxiliary structure that a local change touches
only a local, provably small part of it.

This chapter develops the core of the field. We fix vocabulary
(incremental, decremental, fully dynamic; amortized versus worst-case), then
build the two foundational data structures of the subject: Euler-tour trees for
connectivity in a \emph{forest}, and the Holm--de~Lichtenberg--Thorup
(\textsc{HdLT}) logarithmic-level structure that promotes forest connectivity to
\emph{fully dynamic} connectivity on general graphs in $\Oh(\log^2 n)$ amortized
time. We extend the same machinery to the dynamic minimum spanning tree. We trace
a small graph through every structure by hand. Finally we confront the wall:
conditional lower bounds, organized around the Online Matrix--Vector (OMv)
conjecture, that explain why so many dynamic problems are stuck at roughly
$\sqrt{n}$ or polylog and not better.

\chapterline

\section{A taxonomy of ``dynamic''}
\label{ch39:taxonomy}

\begin{focusbox}
Three update models, in increasing generality.
\begin{itemize}
  \item \textbf{Incremental}: edges are only \emph{inserted} (the graph grows).
  \item \textbf{Decremental}: edges are only \emph{deleted} (the graph shrinks).
  \item \textbf{Fully dynamic}: both insertions and deletions, interleaved
        arbitrarily.
\end{itemize}
Two cost models, with different promises.
\begin{itemize}
  \item \textbf{Amortized}: the bound holds \emph{on average} over a sequence of
        operations; an individual operation may be slow as long as the running
        total is controlled.
  \item \textbf{Worst-case}: the bound holds for \emph{every single} operation,
        no exceptions.
\end{itemize}
\end{focusbox}

The split matters. Incremental connectivity is essentially free: a union-find
(disjoint-set) forest answers ``are $u,v$ connected?'' under edge insertions in
$\Oh(\alpha(n))$ amortized per operation, where $\alpha$ is the inverse Ackermann
function (Chapter~6). Deletions are the hard direction, because union-find has no
un-union; once two components merge, splitting them back apart on an edge deletion
is exactly the operation the structure cannot do. Fully dynamic is hardest of all,
and most of the cleverness in this chapter exists to survive deletions.

\begin{intuitionbox}
Why are deletions so much harder than insertions? An insertion can only
\emph{merge} information -- two components become one, a smaller MST weight
appears -- and merging is monotone and easy to memoize. A deletion can
\emph{shatter} information: cutting one tree edge may or may not disconnect the
graph, depending on whether some \emph{other} edge, currently sitting idle,
happens to reconnect the two halves. To answer that you must, in effect, search
for a replacement edge, and finding that candidate edge without scanning the whole
graph is where the design of dynamic graph algorithms finds its central challenge.
\end{intuitionbox}

\begin{marginnotebox}
Dynamic graph algorithms progressed from $\Oh(\sqrt{m})$ via Frederickson's topology trees (1985) and $\Oh(\sqrt{n})$ via sparsification (1992), to the polylogarithmic era launched by Henzinger and King (1995) and Holm, de~Lichtenberg, and Thorup (1998); randomized worst-case polylog bounds arrived with Kapron, King, and Mountjoy (2013).
\end{marginnotebox}

\section{Warm-up: incremental connectivity and the limits of union-find}
\label{ch39:incremental}

Start where it is easy. We maintain connected components under insertions only.

\begin{algorithm}[H]
\caption{Incremental connectivity via union-find}
\label{ch39:alg-unionfind}
\SetKwFunction{Find}{Find}\SetKwFunction{Union}{Union}\SetKwFunction{Conn}{Connected}
\SetKwProg{Fn}{Function}{:}{}
\Fn{\Union{$u,v$}}{
  $r_u \gets \Find{u}$; \quad $r_v \gets \Find{v}$\;
  \lIf{$r_u = r_v$}{\Return}
  attach the smaller-rank root under the larger (union by rank)\;
}
\Fn{\Conn{$u,v$}}{
  \Return $\Find{u} = \Find{v}$\;
}
\Fn{\Find{$x$}}{
  \If{$\mathrm{par}[x]\neq x$}{
    $\mathrm{par}[x]\gets\Find{$\mathrm{par}[x]$}$ \tcp*[r]{path compression}
  }
  \Return $\mathrm{par}[x]$\;
}
\end{algorithm}

With union by rank and path compression, any sequence of $m$ operations runs in
$\Oh(m\,\alpha(n))$ total, so $\Oh(\alpha(n))$ amortized each. That is effectively
constant. Done -- for insertions.

\begin{trapbox}
The temptation is to ``support deletion'' by storing a \texttt{rank} and undoing
the last union. That works only for a stack-like \emph{rollback} pattern
(union-find with rollback: no path compression, union by rank, undo in
LIFO order, $\Oh(\log n)$ per op). It does \emph{not} give fully dynamic
deletion of an \emph{arbitrary} earlier edge, because the component that edge
helped form may now be held together by many other edges, and removing the wrong
internal pointer corrupts the forest. Rollback union-find is the right tool for
\emph{offline} dynamic connectivity (segment-tree-on-time, below), not for the
online fully dynamic problem.
\end{trapbox}

\begin{connectionbox}
Rollback union-find plus a ``segment tree over the timeline'' solves
\emph{offline} fully dynamic connectivity: each edge is alive on a time interval
$[t_\text{in},t_\text{out})$; place it on the $\Oh(\log q)$ canonical nodes of a
segment tree over time, DFS the segment tree performing unions on entry and
rollbacks on exit, and answer each query at its leaf. Total
$\Oh((m+q)\log q\,\log n)$. This is the dynamic-connectivity analogue of the
offline trick in Chapter~21 (Euler tours) and is worth knowing because online
methods are strictly harder. Online is what the rest of this chapter is about.
\end{connectionbox}

\section{Euler-tour trees: connectivity in a forest}
\label{ch39:ett}

Before we can do general graphs, we need to do \emph{forests} dynamically: support
\textsc{Link}, \textsc{Cut}, and connectivity queries on a forest where the answer
to a query never has to search for a replacement edge (there are no spare edges in
a forest). Chapter~22's link-cut trees already do this, but for connectivity we
prefer a structure that also makes \emph{subtree} aggregates and ``pick any edge
incident to this tree'' easy: the \emph{Euler-tour tree} (ETT) of Henzinger and
King, and Tarjan.

\begin{intuitionbox}
Recall the Euler tour of a tree (Chapter~21): walk the tree depth-first and write
down each vertex every time you touch it. A tree with $n$ vertices has $n-1$
edges; the closed Euler tour visits $2(n-1)$ edges and lists $2n-1$ vertex
occurrences. The key move is to forget the tree shape and keep only this
\emph{sequence}. Then:
\begin{itemize}
  \item ``Are $u$ and $v$ in the same tree?'' $\Leftrightarrow$ ``do their
        occurrences lie in the same sequence?''
  \item \textsc{Link} and \textsc{Cut} become \emph{splits and concatenations} of
        sequences.
\end{itemize}
Store each sequence in a balanced BST keyed by position (a \emph{balanced BST as a
sequence}, i.e. an order-statistic / implicit-key treap or splay tree), and
splits and concatenations are $\Oh(\log n)$.
\end{intuitionbox}

\subsection*{The Euler tour as a sequence}

Represent the Euler tour as a list of \emph{half-edges} and vertex visits. A clean
encoding: for a tree, store the sequence of directed edges of the tour, plus one
self-loop occurrence $\langle v,v\rangle$ per vertex so isolated vertices and
endpoints have an anchor. Each undirected tree edge $\{u,v\}$ appears as two
directed occurrences, $\langle u,v\rangle$ and $\langle v,u\rangle$. We keep, for
each such occurrence, a pointer to its node in the balanced BST.

\begin{figure}[H]
\centering
\begin{tikzpicture}[every node/.style={circle,draw=nodeblue,thick,minimum size=7mm,
   inner sep=0pt,font=\small}, level distance=11mm, sibling distance=15mm]
\node (a) {a}
  child {node (b) {b}}
  child {node (c) {c}
    child {node (d) {d}}};
\begin{scope}[xshift=5.6cm, yshift=-1.0cm]
  \node[draw=none,font=\small\bfseries] at (0,1.4) {Euler tour (rooted at $a$)};
  \node[draw=none,font=\ttfamily\small] at (0,0.6)
    {a\,b\,a\,c\,d\,c\,a};
  \node[draw=none,font=\small] at (0,-0.1)
    {as occurrences: $a_0\,b\,a_1\,c_0\,d\,c_1\,a_2$};
  \node[draw=none,font=\small,align=left] at (0,-1.1)
    {edge $\{a,c\}$ $\to$ occurrences\\ $\langle a,c\rangle$ (enter $c$),
     $\langle c,a\rangle$ (leave $c$)};
\end{scope}
\end{tikzpicture}
\caption{A 4-node tree and one Euler tour. The tour string has $2n-1=7$ vertex
occurrences. Connectivity of $u,v$ is decided by whether their occurrences sit in
the same balanced BST.}
\label{ch39:fig-ett}
\end{figure}

\subsection*{Link and Cut as sequence surgery}

\begin{whybox}
\textbf{Reroot.} It is convenient to be able to make any vertex $u$ the start of
its tour. If the tour reads $S = X\,u\,Y$ where the marked $u$ is some occurrence
of $u$, then the rerooted tour is $u\,Y\,X\,u$. So rerooting at $u$ is: split the
sequence at $u$, swap the two halves, and reattach an occurrence of $u$ at the
end. One split, one concatenate.

\textbf{Link$(u,v)$} (join two different trees by edge $\{u,v\}$). Reroot $u$'s
tour at $u$ and $v$'s tour at $v$. Now $u$'s tour starts with $u$ and $v$'s starts
with $v$. The merged tour is
\[
  \underbrace{u\,\cdots}_{\text{tour}(u)}\; \langle u,v\rangle \;
  \underbrace{v\,\cdots}_{\text{tour}(v)}\; \langle v,u\rangle \; .
\]
Concatenate tour$(u)$, the new occurrence $\langle u,v\rangle$, tour$(v)$, the new
occurrence $\langle v,u\rangle$. A handful of splits/concatenations.

\textbf{Cut$(\{u,v\})$.} The edge contributes occurrences $\langle u,v\rangle$ and
$\langle v,u\rangle$; between them sits exactly the Euler tour of $v$'s subtree.
Splitting the sequence just before $\langle u,v\rangle$ and just after
$\langle v,u\rangle$ isolates that subtour as one piece; concatenating the
outside two pieces removes the edge. Two splits, one concatenate.
\end{whybox}

\begin{algorithm}[H]
\caption{Euler-tour tree operations (balanced-BST-as-sequence)}
\label{ch39:alg-ett}
\SetKwFunction{Reroot}{Reroot}\SetKwFunction{Link}{Link}\SetKwFunction{Cut}{Cut}
\SetKwFunction{Conn}{Connected}\SetKwFunction{Root}{TourRoot}
\SetKwFunction{Split}{Split}\SetKwFunction{Cat}{Concat}
\SetKwProg{Fn}{Function}{:}{}
\Fn{\Reroot{$u$}}{
  let $p$ be the BST node of any occurrence of $u$\;
  $(A,B)\gets \Split$ the sequence just before $p$ \tcp*[r]{$A=X,\;B=u\,Y$}
  \Return $\Cat(B,\,A,\,\text{new occ. }u)$\;
}
\Fn{\Link{$u,v$}}{
  $T_u\gets\Reroot{u}$; \quad $T_v\gets\Reroot{v}$\;
  \Return $\Cat(T_u,\,\langle u,v\rangle,\,T_v,\,\langle v,u\rangle)$\;
}
\Fn{\Cut{$u,v$}}{
  locate occurrences $\langle u,v\rangle$ and $\langle v,u\rangle$\;
  split out the contiguous block from $\langle u,v\rangle$ to $\langle v,u\rangle$ inclusive\;
  concatenate the two outside blocks; discard the two edge occurrences\;
}
\Fn{\Conn{$u,v$}}{
  \Return $\Root{u} = \Root{v}$ \tcp*[r]{root of the BST containing an occurrence}
}
\end{algorithm}

\begin{complexitybox}
Each ETT operation does $\Oh(1)$ splits/concatenations on balanced sequences of
length $\Oh(n)$, hence $\Oh(\log n)$ time (worst-case with a treap/red-black
sequence; amortized with splay trees). Space is $\Oh(n)$: there are $2n-1$
occurrences per tree, summed to $\Oh(n)$ over the forest. \textsc{Connected} is a
``find the root of my BST'' query, $\Oh(\log n)$.
\end{complexitybox}

\begin{labbox}
Trace \textsc{Link}$(b,d)$ on Figure~\ref{ch39:fig-ett} as if $b$ and $d$ were in
\emph{separate} trees (pretend the $c$--$d$ side is its own tree containing $d$).
Reroot tour$(b)=b$ at $b$: still $b$. Reroot the tour $c\,d\,c$ at $d$: it becomes
$d\,c\,d$. Then the merged tour is
$b\,\langle b,d\rangle\,d\,c\,d\,\langle d,b\rangle$. Write out the BST node
sequence and confirm $\textsc{Connected}(b,c)$ now returns true while it was false
before. Doing this once cements why connectivity is just ``same BST.''
\end{labbox}

\begin{aestheticbox}
ETT also maintains \emph{subtree} aggregates directly, because a subtree is a
contiguous range of the tour (the block between $\langle u,v\rangle$ and
$\langle v,u\rangle$). Augment the balanced BST with subtree sums and you get
dynamic subtree-sum on a forest for free. What ETT does \emph{not} give you, and
link-cut trees do, are \emph{root-to-node path} aggregates -- a path is not a
contiguous range of an Euler tour. Pick ETT for subtree/connectivity, LCT for
path. This is the cleanest way to remember which dynamic-tree structure to reach
for.
\end{aestheticbox}

\section{Fully dynamic connectivity: the HdLT logarithmic levels}
\label{ch39:hdlt}

Now the centerpiece. We want fully dynamic connectivity on a \emph{general} graph:
insert edge, delete edge, query ``connected?'' We maintain a spanning forest $F$
of the current graph. An insertion is easy. The hard case is deleting a
\emph{tree} edge of $F$: we must decide whether some non-tree edge can replace it,
and if so install that replacement, all without scanning everything.

\begin{intuitionbox}
Keep a spanning forest $F$ inside the graph $G$. Non-tree edges are
``\emph{spares}'': they connect two vertices already connected through $F$.
\begin{itemize}
  \item \textbf{Insert $\{u,v\}$}: if $u,v$ already connected in $F$, store
        $\{u,v\}$ as a spare; else \textsc{Link} it into $F$ as a tree edge.
  \item \textbf{Query}: connectivity in $F$ (an ETT query).
  \item \textbf{Delete a spare}: just remove it.
  \item \textbf{Delete a tree edge $\{u,v\}$}: \textsc{Cut} it, splitting $F$'s
        tree into parts $T_u\ni u$ and $T_v\ni v$. Search for a spare edge with
        one endpoint in $T_u$ and the other in $T_v$. If found, promote it to a
        tree edge (\textsc{Link}); if none exists, the deletion genuinely
        disconnected the graph.
\end{itemize}
Everything hinges on finding that replacement spare cheaply. A naive search scans
all spares incident to the smaller side, which can be $\Th(m)$. The HdLT idea is a
charging scheme that guarantees each spare is scanned a bounded number of times
over its whole life.
\end{intuitionbox}

\subsection*{Levels and the smaller-half rule}

\begin{focusbox}
Assign every edge a \textbf{level} $\ell(e)\in\{0,1,\dots,L\}$ with
$L=\floor{\log_2 n}$. Maintain, for each level $i$, a spanning forest $F_i$ of the
subgraph $G_i$ formed by edges of level $\geq i$. Two invariants:
\begin{enumerate}
  \item[(I1)] $F_0\supseteq F_1\supseteq\cdots\supseteq F_L$; each $F_i$ is a
        spanning forest of $G_i$, and $F:=F_0$ is a spanning forest of all of $G$.
  \item[(I2)] Every connected component of the subgraph induced by level-$\geq i$
        edges has at most $\floor{n/2^i}$ vertices.
\end{enumerate}
A tree edge of $F$ at level $i$ is also a tree edge of $F_i,F_{i-1},\dots,F_0$.
We keep an ETT for each level's forest $F_i$, augmented so each tour node knows
whether its vertex has any incident \emph{non-tree} edge at that level.
\end{focusbox}

Invariant (I2) is the engine (Figure~\ref{ch39:fig-levels}). Because a level-$i$
component has $\leq n/2^i$ vertices and levels go up to $\log n$, an edge can be
\emph{promoted} (its level raised) at most $\log n$ times before its component
would be too small to hold it.
Each promotion is what pays for the work of scanning it. That is the whole
amortized argument in one sentence; the rest is bookkeeping.

\begin{figure}[H]
\centering
\begin{tikzpicture}[
  scale=0.92,
  vnode/.style={circle,draw=sourceblue!80!black,fill=sourceblue!15,minimum size=5.5mm,inner sep=0pt,font=\footnotesize\bfseries},
  lanehead/.style={font=\footnotesize\scshape,anchor=west},
  cutmark/.style={trapred,line width=1.5pt}
]
  % Lane 0 (Top)
  \fill[sourceblue!6,rounded corners=4pt,draw=sourceblue!30] (-0.8, 2.5) rectangle (12.2, 5.95);
  \node[lanehead,text=sourceblue!90!black] at (-0.6, 5.6) {Level $\ell=0$ ($F_0 \subseteq G_0$, component size $\le n$)};
  
  % Smaller component T_u
  \draw[quotientorange!30,fill=quotientorange!12,rounded corners=4pt,dashed] (0.3, 2.65) rectangle (3.5, 4.35);
  \node[font=\tiny\scshape,text=quotientorange!90!black] at (1.9, 2.85) {Smaller Side $T_u$ ($|T_u| \le n/2$)};
  \node[vnode] (u1) at (0.8, 3.85) {$1$};
  \node[vnode] (u2) at (3.0, 3.85) {$2$};
  \draw[sourceblue,line width=1.5pt] (u1) -- node[above=1pt,font=\tiny]{tree} (u2);
  
  % Larger component T_v
  \draw[gray!30,fill=gray!6,rounded corners=4pt] (4.8, 2.95) rectangle (11.5, 4.55);
  \node[font=\tiny\scshape,text=gray!70!black] at (8.15, 4.3) {Larger Side $T_v$};
  \node[vnode] (v3) at (5.5, 3.85) {$3$};
  \node[vnode] (v4) at (7.2, 3.85) {$4$};
  \node[vnode] (v5) at (9.0, 3.85) {$5$};
  \node[vnode] (v6) at (10.7, 3.85) {$6$};
  \draw[sourceblue,line width=1.5pt] (v3) -- (v4);
  \draw[sourceblue,line width=1.5pt] (v4) -- (v5);
  \draw[sourceblue,line width=1.5pt] (v5) -- (v6);
  
  % Cut tree edge (2, 3)
  \draw[sourceblue!40,line width=1.2pt,dashed] (u2) -- (v3);
  \draw[cutmark] (3.9, 4.1) -- (4.4, 3.6);
  \draw[cutmark] (3.9, 3.6) -- (4.4, 4.1);
  \node[font=\tiny\bfseries,text=trapred,below=2pt] at (4.15, 3.55) {\textsc{Cut} $\{2,3\}$};

  % Crossing replacement edge (1, 4)
  \draw[targetgreen!80!black,line width=1.8pt,bend left=22] (u1.north east) to node[above=2pt,font=\tiny\bfseries,text=targetgreen!90!black]{replacement tree edge} (v4.north west);

  % Internal non-tree candidate edge (1, 2) spare
  \draw[quotientorange!80!black,thick,dashed,bend right=24] (u1.south east) to node[below=1pt,font=\tiny,text=quotientorange!90!black]{internal $\to$ promote} (u2.south west);

  % Arrow showing promotion to Level 1
  \draw[->,quotientorange!90!black,very thick] (1.9, 2.6) -- node[right,font=\tiny\bfseries]{promote internal edges} (1.9, 2.1);

  % Lane 1 (Middle)
  \fill[quotientorange!6,rounded corners=4pt,draw=quotientorange!30] (-0.8, -0.1) rectangle (12.2, 2.1);
  \node[lanehead,text=quotientorange!90!black] at (-0.6, 1.8) {Level $\ell=1$ ($F_1 \subseteq G_1$, components $\le n/2$)};
  \draw[quotientorange!40,fill=quotientorange!15,rounded corners=4pt] (0.3, 0.2) rectangle (3.5, 1.5);
  \node[vnode] (lu1) at (0.8, 0.95) {$1$};
  \node[vnode] (lu2) at (3.0, 0.95) {$2$};
  \draw[quotientorange!90!black,line width=1.5pt] (lu1) -- node[below=1pt,font=\tiny]{promoted tree} (lu2);
  \node[font=\tiny,text=gray!70!black] at (7.5, 0.7) {Edges touching $T_u$ that failed to cross enter $F_1$ (safe since $|T_u| \le n/2$)};

  % Lane 2 (Bottom)
  \fill[targetgreen!6,rounded corners=4pt,draw=targetgreen!30] (-0.8, -1.8) rectangle (12.2, -0.5);
  \node[lanehead,text=targetgreen!90!black] at (-0.6, -0.8) {Level $\ell=2 \dots \floor{\log_2 n}$ (components $\le n/2^\ell$)};
  \node[font=\tiny,text=gray!70!black] at (5.7, -1.3) {Climbing stops at level $\floor{\log_2 n}$: each edge promoted at most $\Oh(\log n)$ times};
\end{tikzpicture}
\caption{The Holm--de~Lichtenberg--Thorup (HdLT) logarithmic level hierarchy. When tree edge $\{2,3\}$ is cut at level~$0$, the smaller component $T_u$ is inspected. Candidate edges with both endpoints in $T_u$ are promoted to level~$1$ (invariant (I2) holds because $|T_u| \le n/2$). A candidate edge $\{1,4\}$ crossing between $T_u$ and $T_v$ is immediately selected as the replacement tree edge, reconnecting $F_0$. Because an edge climbs at most $\floor{\log_2 n}$ levels, every edge scan is charged to a promotion, yielding $\Oh(\log^2 n)$ amortized time.}
\label{ch39:fig-levels}
\end{figure}

\begin{algorithm}[H]
\caption{HdLT fully dynamic connectivity}
\label{ch39:alg-hdlt}
\SetKwFunction{Ins}{Insert}\SetKwFunction{Del}{Delete}\SetKwFunction{Conn}{Connected}
\SetKwFunction{Replace}{Replace}\SetKwFunction{Link}{Link}\SetKwFunction{Cut}{Cut}
\SetKwProg{Fn}{Function}{:}{}
\Fn{\Conn{$u,v$}}{\Return $u,v$ in same tree of $F_0$ (ETT query at level $0$)\;}
\Fn{\Ins{$u,v$}}{
  $\ell(\{u,v\})\gets 0$\;
  \lIf{\Conn{$u,v$}}{store $\{u,v\}$ as a non-tree edge at level $0$}
  \lElse{add $\{u,v\}$ to $F_0$ as a tree edge (ETT \Link\ at level $0$)}
}
\Fn{\Del{$u,v$}}{
  $i\gets \ell(\{u,v\})$\;
  \uIf{$\{u,v\}$ is a non-tree edge}{remove it from the level-$i$ adjacency\;}
  \Else{
    remove $\{u,v\}$ from $F_i,F_{i-1},\dots,F_0$ (ETT \Cut\ at each level)\;
    \Replace{$u,v,i$}\;
  }
}
\Fn{\Replace{$u,v,i$}}{
  \For{$j\gets i$ \textbf{\emph{downto}} $0$}{
    let $T_u,T_v$ be the two pieces of the level-$j$ forest split by the cut\;
    let $T$ be the \emph{smaller} of $T_u,T_v$ (fewer vertices, via ETT subtree size)\;
    \ForEach{level-$j$ tree edge $e$ inside $T$}{
      raise $\ell(e)$ to $j{+}1$; add $e$ to $F_{j+1}$ \tcp*[r]{promote tree edges}
    }
    \ForEach{level-$j$ non-tree edge $e=\{x,y\}$ with $x\in T$}{
      \uIf{$y$ on the other side}{
        add $e$ to $F_j,\dots,F_0$ as the replacement tree edge\;
        \Return \tcp*[r]{reconnected}
      }
      \Else{raise $\ell(e)$ to $j{+}1$ \tcp*[r]{both endpoints in $T$: promote}}
    }
  }
  \tcp{no replacement at any level: $u,v$ are now genuinely disconnected}
}
\end{algorithm}

\subsection*{Why \textsc{Replace} is correct}

\begin{whybox}
When we cut a level-$i$ tree edge, we look for a replacement starting at the
\emph{highest} level $i$ and descend. At level $j$ we have split the level-$j$
forest into $T_u,T_v$. We scan from the smaller side $T$. Two things can happen to
each scanned edge:
\begin{itemize}
  \item It is a non-tree edge crossing to the other side: it reconnects the two
        pieces. Use it, at level $j$, and stop. Correct because any spare crossing
        $T_u\!:\!T_v$ at this level restores connectivity, and a crossing edge of
        the \emph{highest} possible level is safe to install without violating
        (I1)/(I2) below it.
  \item It is fully inside $T$ (a tree edge of $T$, or a non-tree edge with both
        endpoints in $T$): we \emph{promote} it to level $j{+}1$. This is allowed
        because $T$ is the smaller half, so after the cut $\abs{T}\leq
        \floor{\tfrac12\abs{T_u\cup T_v}}\leq \floor{\tfrac12\floor{n/2^j}} =
        \floor{n/2^{j+1}}$, which is exactly the size budget (I2) demands at level
        $j{+}1$ (the last step is the integer identity
        $\floor{\floor{x}/2}=\floor{x/2}$). Promotion never breaks (I2).
\end{itemize}
If level $j$ yields no crossing spare, every level-$j$ spare touching $T$ has been
promoted, so none crosses; descend to $j{-}1$ and repeat on the coarser forest.
If we exhaust level $0$ without a crossing edge, then no edge of the graph crosses
$T_u\!:\!T_v$ at all -- the deletion truly disconnected the graph.
\end{whybox}

\begin{formalbox}
\textbf{Theorem (HdLT).} Fully dynamic connectivity is maintainable in
$\Oh(\log^2 n)$ amortized time per update and $\Oh(\log n /\log\log n)$ time per
query, using $\Oh(m+n\log n)$ space.

\textbf{Cost of a promotion.} Each edge starts at level $0$ and can be promoted at
most $L=\Oh(\log n)$ times, by (I2). Charge the cost of \emph{scanning} an edge
during \textsc{Replace} to a promotion: every edge we scan and keep is promoted
one level (the only scanned edge we do \emph{not} promote is the single crossing
replacement, which we charge to the deletion that caused it). So the total number
of edge-scan events over the whole run is $\Oh(m\log n)$ scans, each costing an
ETT update of $\Oh(\log n)$ to move the edge between levels. That is
$\Oh(m\log^2 n)$ total promotion work, i.e. $\Oh(\log^2 n)$ amortized per
operation. Insertions and the $L$-fold \textsc{Cut}/\textsc{Link} along the level
forests add only $\Oh(\log^2 n)$ each. The query bound uses a slightly cleverer
ETT root-marking; the simple version gives $\Oh(\log n)$.
\end{formalbox}

\begin{whybox}
\textbf{Finding a vertex of $T$ that still has a level-$j$ non-tree edge.}
We do not literally iterate all vertices of $T$ to find incident spares. Augment
each level-$j$ ETT so every BST subtree records ``does my range contain an
occurrence whose vertex has a level-$j$ non-tree edge?'' Then ``find any non-tree
edge incident to $T$'' is a single $\Oh(\log n)$ descent into the BST of $T$. We
pop edges from a per-vertex adjacency list one at a time, which is what lets us
charge each examined edge to its promotion. This augmentation is the bridge
between ``ETT supports subtree aggregates'' and ``HdLT finds replacement edges
fast.''
\end{whybox}

\subsection*{Worked example}

\begin{figure}[H]
\centering
\begin{tikzpicture}[every node/.style={circle,draw=nodeblue,thick,minimum size=7mm,
   inner sep=0pt,font=\small}, node distance=14mm]
\node (a) {a}; \node (b) [right=of a] {b}; \node (c) [right=of b] {c};
\node (d) [below=of a] {d}; \node (e) [below=of b] {e}; \node (f) [below=of c] {f};
% tree edges (level 0, solid thick)
\draw[very thick,nodeblue] (a)--(b); \draw[very thick,nodeblue] (b)--(c);
\draw[very thick,nodeblue] (a)--(d); \draw[very thick,nodeblue] (b)--(e);
\draw[very thick,nodeblue] (c)--(f);
% spare edges (dashed)
\draw[dashed,trapred,thick] (d)--(e); \draw[dashed,trapred,thick] (e)--(f);
\node[draw=none,font=\small,trapred] at (1.7,-3.2) {dashed = non-tree (spare) edges};
\node[draw=none,font=\small,nodeblue] at (1.7,-3.8) {thick blue = spanning forest $F_0$};
\end{tikzpicture}
\caption{A graph on $\{a,\dots,f\}$. Spanning forest $F_0$ is the five blue tree
edges; $\{d,e\}$ and $\{e,f\}$ are spares (level $0$).}
\label{ch39:fig-hdlt}
\end{figure}

\begin{labbox}
Delete the tree edge $\{a,b\}$ in Figure~\ref{ch39:fig-hdlt}. It is a level-$0$
tree edge, so \textsc{Cut} it; $F_0$ splits into $T_a=\{a,d\}$ and
$T_b=\{b,c,e,f\}$. The smaller side is $T_a$ (2 vertices). Scan level-$0$ edges
touching $T_a$:
\begin{itemize}
  \item tree edge $\{a,d\}$: inside $T_a$, promote it to level $1$;
  \item spare $\{d,e\}$: endpoint $e\in T_b$ -- \emph{crossing}. Install it as the
        replacement tree edge.
\end{itemize}
$F_0$ becomes $\{a,d\},\{d,e\},\{b,e\},\{b,c\},\{c,f\}$, one connected tree again:
the deletion did \emph{not} disconnect the graph, and we found the replacement
after touching just two edges from the smaller side. Now \emph{also} delete
$\{d,e\}$ and $\{e,f\}$ and then $\{a,b\}$ again, and watch a deletion finally
report a real disconnection. Tracing this by hand is the fastest way to internalize
the smaller-half charging.
\end{labbox}

\begin{trapbox}
Three places people get HdLT wrong. (1) Scanning the \emph{larger} half: the
size-halving budget (I2) only holds for the smaller half; scan the wrong side and
promotions can violate the level cap, destroying the amortized bound. (2)
Forgetting to promote the level-$j$ \emph{tree} edges of $T$ alongside the
spares -- if tree edges are left behind, $F_{j+1}$ stops being a spanning forest
of $G_{j+1}$ and (I1) breaks. (3) Installing a replacement at the wrong level: the
replacement edge must enter $F_j,\dots,F_0$ (all levels at and below where it was
found), or the lower forests lose connectivity.
\end{trapbox}

\begin{marginnotebox}
The smaller-half promotion charging mirrors union-by-size and small-to-large merging. HdLT combined this charging scheme with Euler-tour trees so that finding a replacement edge becomes a logarithmic-time BST descent, reducing the update bound from $\Oh(\sqrt n)$ to $\Oh(\log^2 n)$.
\end{marginnotebox}

\section{Dynamic minimum spanning tree}
\label{ch39:dynamic-mst}

Maintaining a minimum spanning forest under insertions and deletions is the same
skeleton with one extra idea: replacement edges must be chosen by \emph{weight}.

\begin{intuitionbox}
Two classical facts (Chapter~6) drive everything.
\begin{itemize}
  \item \textbf{Cycle rule.} Inserting edge $e$ creates a cycle in the current MST
        $T$ iff its endpoints are already connected. If so, $e$ replaces the
        \emph{maximum-weight} edge on the tree path between its endpoints, but only
        when $e$ is lighter than that maximum; otherwise $e$ becomes a spare.
  \item \textbf{Cut rule.} Deleting a tree edge $e$ splits $T$ into two sides; the
        correct replacement is the \emph{minimum-weight} spare crossing the cut.
\end{itemize}
\end{intuitionbox}

Figure~\ref{ch39:fig-dmst-rules} contrasts the two dual update primitives side by side: path-max replacement on an edge insertion versus lightest-crossing replacement on a tree-edge deletion.

\begin{figure}[H]
\centering
\begin{tikzpicture}[
  vx/.style={circle,draw=nodeblue,fill=nodeblue!12,thick,minimum size=6.5mm,inner sep=1pt,font=\scriptsize\bfseries}]

  % Panel 1: Cycle Rule (Insertion)
  \begin{scope}[shift={(0,0)}]
    \draw[rounded corners=6pt,fill=sourceblue!4,draw=sourceblue!25,thick] (-0.8,-1.85) rectangle (4.8,1.55);
    \node[font=\footnotesize\bfseries,text=sourceblue] at (2.0,1.85) {1.\ Cycle Rule: Insert $e=\{u,v\}$ ($w=3$)};

    \node[vx] (u1) at (0,0) {$u$};
    \node[vx] (x1) at (0.95,0.9) {$x$};
    \node[vx] (y1) at (3.05,0.9) {$y$};
    \node[vx] (v1) at (4.0,0) {$v$};

    \draw[sourceblue,very thick] (u1)--node[above left=-1pt,font=\tiny]{$2$}(x1);
    \draw[trapred,very thick,dashed] (x1)--node[above=2pt,font=\tiny\bfseries]{$e^\star=7$ (cut)}(y1);
    \draw[sourceblue,very thick] (y1)--node[above right=-1pt,font=\tiny]{$4$}(v1);
    \draw[targetgreen,very thick] (u1)--node[below=2pt,font=\tiny\bfseries]{insert $e=3 < 7$}(v1);
    \node[font=\scriptsize,text=edgegray] at (2.0,-1.45) {LCT \texttt{path-max} finds $e^\star=7$; swap $3 \leftrightarrow 7$};
  \end{scope}

  % Panel 2: Cut Rule (Deletion)
  \begin{scope}[shift={(6.4,0)}]
    \draw[rounded corners=6pt,fill=quotientorange!4,draw=quotientorange!25,thick] (-0.8,-1.85) rectangle (4.8,1.55);
    \node[font=\footnotesize\bfseries,text=quotientorange] at (2.0,1.85) {2.\ Cut Rule: Delete Tree Edge ($w=2$)};

    \node[vx] (u2) at (0,0.8) {$u$};
    \node[vx] (d2) at (0,-0.5) {$d$};
    \node[vx] (v2) at (4.0,0.8) {$v$};
    \node[vx] (f2) at (4.0,-0.5) {$f$};

    \draw[sourceblue,very thick] (u2)--(d2) (v2)--(f2);
    \draw[trapred,thick,dashed] (u2)--node[above=1pt,font=\tiny\bfseries]{deleted ($2$)}(v2);
    \draw[targetgreen,very thick] (d2)--node[sloped,above=1pt,font=\tiny\bfseries]{lightest spare $6$}(v2);
    \draw[edgegray,dashed,thick] (d2)--node[below=1pt,font=\tiny]{spare $9$}(f2);
    \node[font=\scriptsize,text=edgegray] at (2.0,-1.45) {leveled ETT promotes $6$ across $T_u : T_v$};
  \end{scope}
\end{tikzpicture}
\caption{Dynamic MST update rules. Left (Cycle Rule): Inserting $e=\{u,v\}$ of weight $3$ closes a cycle with the tree path $u\text{--}x\text{--}y\text{--}v$; a link-cut tree \texttt{path-max} query locates the heaviest path edge $e^\star=\{x,y\}$ (weight $7>3$) and replaces it with $e$. Right (Cut Rule): Deleting tree edge $\{u,v\}$ splits the tree into $T_u$ and $T_v$; the leveled replacement search installs the lightest crossing spare (weight $6 < 9$).}
\label{ch39:fig-dmst-rules}
\end{figure}

\subsection*{Insertion: the cycle rule with path-max}

\begin{algorithm}[H]
\caption{Dynamic MST -- edge insertion}
\label{ch39:alg-dmst-ins}
\SetKwFunction{InsMST}{InsertMST}\SetKwFunction{Link}{Link}\SetKwFunction{Cut}{Cut}
\SetKwProg{Fn}{Function}{:}{}
\Fn{\InsMST{$u,v,w$}}{
  \uIf{$u,v$ in different trees of $T$}{\Link{$u,v$} with weight $w$ \tcp*[r]{forest grows}}
  \Else{
    $e^\star\gets$ maximum-weight edge on the $u$--$v$ tree path \tcp*[r]{$\Oh(\log n)$ via LCT path-max}
    \uIf{$w < \mathrm{weight}(e^\star)$}{
      \Cut{$e^\star$}; demote $e^\star$ to a spare; \Link{$u,v$} with weight $w$\;
    }
    \lElse{store $\{u,v\}$ as a spare of weight $w$}
  }
}
\end{algorithm}

The path-max query and the cut/link are exactly the link-cut tree primitives of
Chapter~22: a link-cut tree stores the represented MST and answers ``heaviest edge
on the path $u$--$v$'' in $\Oh(\log n)$ amortized. So \emph{incremental} MST is
$\Oh(\log n)$ per insertion -- clean, and worth contrasting with the deletion
case, which is the hard one.

\subsection*{Deletion: weighted replacement via leveled HdLT}

Deleting a tree edge needs the \emph{minimum-weight} crossing spare, which the
plain HdLT replacement search does not order by weight. The fix (Holm--de
Lichtenberg--Thorup, same paper) keeps the leveled-forest structure but, when
searching for a replacement at a given level, takes the lightest crossing spare
found. To make ``lightest crossing'' efficient, the levels are organized so that
the search examines candidate spares in nondecreasing weight, and a careful
argument shows the amortized cost remains $\Oh(\log^2 n)$ per update for the
minimum spanning \emph{forest}.

\begin{focusbox}
\textbf{Dynamic MSF result.} The minimum spanning forest of a fully dynamic graph
is maintainable in $\Oh(\log^2 n)$ amortized time per insertion \emph{and} per
deletion (with the same leveled-ETT machinery, replacement edges chosen by
weight). The MST \emph{weight} is then available in $\Oh(1)$ by keeping a running
sum that is updated on every \textsc{Link}/\textsc{Cut}.
\end{focusbox}

\begin{whybox}
\textbf{Why weight-ordered replacement is correct.} After cutting tree edge $e$,
the two sides are $T_u,T_v$. By the cut rule, the unique correct replacement is the
minimum-weight edge of $G$ crossing $T_u\!:\!T_v$. Any \emph{tree} edge of the old
$T$ other than $e$ stays in the MST (removing $e$ does not change which side those
edges connect), so the only candidates are spares plus possibly nothing. Taking
the lightest crossing spare therefore reproduces what Kruskal/Prim would pick on
the reduced graph. The leveled structure only changes \emph{how} we find that
lightest crosser efficiently; correctness is the cut rule, full stop.
\end{whybox}

\begin{labbox}
On Figure~\ref{ch39:fig-hdlt}, attach weights:
$ab{=}1,bc{=}2,ad{=}3,be{=}4,cf{=}5$ (tree), $de{=}6,ef{=}7$ (spares). Delete
tree edge $be$ (weight 4). Sides: the part containing $e$ versus the rest.
Crossing spares: $de$ (6) and $ef$ (7). The cut rule installs the lighter,
$de=6$. New MSF weight $=1+2+3+5+6=17$. Now instead delete $ab$ (weight 1).
Crossing spares across $\{a,d\}\!:\!\{b,c,e,f\}$: only $de=6$. Replace; weight
becomes $2+3+4+5+6=20$. Verify by recomputing the MST from scratch each time --
they must match, and they do.
\end{labbox}

\begin{connectionbox}
Dynamic MST is where Chapter~22 (link-cut trees, for the cheap incremental side
and path-max) and this chapter's leveled ETT (for the hard decremental side) meet.
The same path-max primitive powers Sleator--Tarjan dynamic-tree speedups of
blocking-flow max-flow (Chapter~12--13). And the offline ``segment tree over the
timeline'' from \S\ref{ch39:incremental}, combined with rollback union-find that
tracks the maximum edge on each merge, gives an elegant offline dynamic-MST that
is often the practical choice when all updates are known in advance.
\end{connectionbox}

\section{Sparsification: trading \texorpdfstring{$m$}{m} for \texorpdfstring{$n$}{n}}
\label{ch39:sparsification}

Many bounds above are stated in $n$, yet a graph can have $m=\Th(n^2)$ edges. The
reason updates do not degrade to $m$ is \emph{sparsification}, the
Eppstein--Galil--Italiano--Nissenzweig technique that deserves a paragraph because
it is the general-purpose lever.

\begin{intuitionbox}
For a property decided by a spanning forest (connectivity, MSF, bipartiteness,
$k$-edge-connectivity for small $k$), the answer depends only on $\Oh(n)$
\emph{certificate} edges, not all $m$. Build a balanced binary tree over the edge
set. Each leaf holds $\Oh(n)$ edges; each internal node stores a \emph{sparse
certificate} of the union of its children's certificates -- itself only $\Oh(n)$
edges. An edge update touches one leaf and the $\Oh(\log m)$ certificates on the
path to the root, recomputing each from two $\Oh(n)$-edge inputs. The dynamic
structure then runs on a graph that always has $\Oh(n)$ edges, so an
$\Oh(\sqrt m)$ per-operation method becomes $\Oh(\sqrt n)$, and $m$ disappears from
polylog bounds entirely.
\end{intuitionbox}

\begin{complexitybox}
Sparsification multiplies the per-operation cost by $\Oh(\log m)$ but replaces
every $m$ inside the base structure by $n$. It is the reason Frederickson's
$\Oh(\sqrt m)$ topology trees became the $\Oh(\sqrt n)$ worst-case bound, and a
clean way to keep $m$ out of the polylog connectivity bounds even on dense graphs.
\end{complexitybox}

\section{Worst-case versus amortized, and a glance at harder problems}
\label{ch39:worstcase}

HdLT is \emph{amortized}: a single deletion can trigger an avalanche of promotions
and run in $\Th(n)$ time, even though the average is $\Oh(\log^2 n)$. For
real-time systems (an individual operation must never stall) that is unacceptable,
and worst-case bounds are genuinely harder to obtain.

\begin{focusbox}
\textbf{Complexity bounds for fully dynamic connectivity.}
\begin{itemize}
  \item Amortized: $\Oh(\log^2 n)$ (HdLT, deterministic), improved to
        $\Oh(\log^2 n/\log\log n)$ (Wulff-Nilsen) and below with refinements.
  \item Worst-case: $\Oh(\sqrt n)$ deterministic (Frederickson + sparsification);
        $\Oh(\poly\log n)$ worst-case \emph{randomized} via cut-set / XOR
        sketching (Kapron--King--Mountjoy 2013 and successors).
  \item Lower bound: $\Om(\log n)$ per operation (Patrascu--Demaine), so no
        dynamic-connectivity structure can be truly constant-time.
\end{itemize}
\end{focusbox}

\begin{aestheticbox}
\textbf{The cut-set sketch}, in one idea. Give each edge $\{u,v\}$ a random
$\Oh(\log n)$-bit id, and store at each vertex the XOR of ids of its incident
edges. For a vertex set $S$, the XOR of the stored values over $S$ cancels every
edge with both endpoints inside $S$ (each such id appears twice) and leaves the
XOR of edges crossing the cut $(S,\bar S)$. If exactly one edge crosses, that XOR
\emph{is} its id -- you have recovered a crossing edge with no search at all.
Layer $\Oh(\log n)$ independent samplings to handle ``more than one crosser,'' and
you get a replacement edge in worst-case polylog time, with high probability. This
is the seed of worst-case dynamic connectivity and a beautiful instance of
linear-sketching meeting graphs (Chapter~40).
\end{aestheticbox}

\section{Conditional lower bounds: the OMv wall}
\label{ch39:omv}

Why do so many dynamic problems plateau at $\sqrt n$ or polylog and refuse to go
faster despite decades of effort? Because, conjecturally, they cannot. The cleanest
explanation is the \emph{Online Matrix--Vector multiplication} (OMv) conjecture of
Henzinger, Krinninger, Nanongkai and Saranurak (2015).

\begin{focusbox}
\textbf{OMv problem.} Fix an $n\times n$ Boolean matrix $M$ (preprocessing
allowed). Then receive vectors $v_1,\dots,v_n\in\{0,1\}^n$ \emph{online}: you must
output $Mv_t$ (Boolean matrix--vector product, $\vee/\wedge$) before seeing
$v_{t+1}$.

\textbf{OMv conjecture.} No algorithm solves OMv in total time
$\Oh(n^{3-\eps})$ for any constant $\eps>0$. (Even with arbitrary polynomial
preprocessing of $M$.)
\end{focusbox}

The naive algorithm computes each $Mv_t$ in $\Oh(n^2)$ for a total of $\Oh(n^3)$;
the conjecture says you cannot shave a polynomial factor off that when the queries
are online. Why believe it? Because offline (all $v_t$ known at once) it
\emph{is} beatable via fast matrix multiplication, yet decades of effort have not
produced a truly subcubic \emph{online} algorithm -- the online constraint seems
to forbid exactly the batching that fast matrix multiplication needs.

\begin{whybox}
\textbf{From OMv to a dynamic lower bound: the reduction pattern}
(Figure~\ref{ch39:fig-omv}). Encode the
matrix $M$ as a directed bipartite graph $H$ on parts $A=\{a_1,\dots,a_n\}$,
$B=\{b_1,\dots,b_n\}$ with a directed edge $b_j\!\to\!a_i$ iff $M_{ij}=1$. To
answer an online query vector $v_t$:
\begin{enumerate}
  \item \emph{Update phase}: for each $j$ with $(v_t)_j=1$, insert a directed edge
        from a fresh source $s$ to $b_j$ (``activate'' column $j$). That is
        $\Oh(n)$ edge insertions.
  \item \emph{Query phase}: for each $i$, a single dynamic query ``is $a_i$
        reachable from $s$?'' reveals whether $\bigvee_j (M_{ij}\wedge (v_t)_j)$,
        i.e. the $i$-th entry of $Mv_t$. That is $\Oh(n)$ queries.
  \item \emph{Cleanup}: delete the $\Oh(n)$ inserted $s$-edges to reset for
        $v_{t+1}$.
\end{enumerate}
So one online matrix--vector product becomes $\Oh(n)$ updates and $\Oh(n)$ queries.
If a dynamic reachability structure had $u(n)$ update and $q(n)$ query time, the
whole OMv instance (with $n$ query vectors) would run in
$\Oh(n\cdot n\cdot(u(n)+q(n)))=\Oh(n^2(u+q))$. For this to beat $n^{3-\eps}$ we
would need $u(n)+q(n)=\Oh(n^{1-\eps})$. Hence: \emph{under OMv, fully dynamic
$s$--$t$ reachability cannot have both update and query in $\Oh(n^{1-\eps})$.}
\end{whybox}

\begin{figure}[H]
\centering
\begin{tikzpicture}[every node/.style={circle,draw=nodeblue,thick,minimum size=7mm,
   inner sep=0pt,font=\small},
   >={Stealth[length=2mm]}, node distance=12mm]
% source
\node (s) at (-2.4,-1.5) {$s$};
% B column (activated by v_t)
\node (b1) at (0,0) {$b_1$};
\node (b2) at (0,-1.5) {$b_2$};
\node (b3) at (0,-3.0) {$b_3$};
% A column
\node (a1) at (3,0) {$a_1$};
\node (a2) at (3,-1.5) {$a_2$};
\node (a3) at (3,-3.0) {$a_3$};
% activation edges (v_t = (1,0,1)): s->b1, s->b3
\draw[->,trapred,thick] (s)--(b1);
\draw[->,trapred,thick] (s)--(b3);
% matrix edges b_j -> a_i for M_{ij}=1
\draw[->,nodeblue] (b1)--(a1);
\draw[->,nodeblue] (b1)--(a2);
\draw[->,nodeblue] (b3)--(a2);
\draw[->,nodeblue] (b2)--(a3);
\node[draw=none,font=\small,trapred,anchor=west] at (-2.0,-3.75)
  {red: $s\!\to\!b_j$ when $(v_t)_j{=}1$};
\node[draw=none,font=\small,nodeblue,anchor=west] at (4.0,-1.5)
  {blue: $b_j\!\to\!a_i$ when $M_{ij}{=}1$};
\end{tikzpicture}
\caption{The OMv gadget for $v_t=(1,0,1)$. A directed query ``is $a_i$ reachable
from $s$?'' returns the $i$-th bit of $Mv_t$: here $a_1$ and $a_2$ are reachable
(via $b_1$, and via $b_1$ or $b_3$), while $a_3$ is not, since its only in-edge is
from the inactive $b_2$. So $Mv_t=(1,1,0)$.}
\label{ch39:fig-omv}
\end{figure}

\begin{formalbox}
\textbf{Consequences (all conditional on OMv).} There is no fully dynamic
algorithm with $\Oh(n^{1-\eps})$ per operation (for constant $\eps>0$) for, among
others:
\begin{itemize}
  \item single-source / $s$--$t$ \textbf{reachability} in directed graphs;
  \item \textbf{strong connectivity} (is the whole digraph one SCC?);
  \item maintaining whether the graph has a \textbf{perfect matching};
  \item $\Oh(\poly\log)$-approximate shortest paths and many distance problems;
  \item subgraph-detection / triangle-type counting under updates.
\end{itemize}
The $\eps$ matters: these bounds rule out genuinely sublinear-per-operation
algorithms, which is exactly why $\Oh(\sqrt n)$ (for the undirected, easier
problems that escape the reduction) or $\Oh(n)$ (for the directed ones it traps)
are the natural ceilings people hit.
\end{formalbox}

\begin{intuitionbox}
The moral. \emph{Undirected} connectivity escapes OMv -- the reduction needs the
asymmetry of reachability, and undirected connectivity's spanning-forest structure
is too well-behaved to encode matrix--vector products -- which is precisely why we
\emph{can} do it in polylog. \emph{Directed} reachability, exact shortest paths,
and matching get caught in the OMv net, which is why their dynamic versions are
stuck near $\Oh(n)$ per update. When a dynamic problem refuses to go below
$\sqrt n$ or $n$ after serious effort, suspect an OMv (or the related 3SUM / APSP)
hardness reduction lurking underneath.
\end{intuitionbox}

\begin{marginnotebox}
Conditional lower bounds based on the OMv conjecture unified dynamic hardness around 2015, alongside SETH, 3SUM, and APSP, proving that problems stuck at $\sqrt{n}$ or linear update times share common computational barriers under fine-grained complexity.
\end{marginnotebox}

\begin{connectionbox}
The OMv barrier is the dynamic-graph face of \emph{fine-grained complexity}
(Chapter~37 covers NP-hardness; this is the polynomial-time analogue). It also
explains the design of \S\ref{ch39:sparsification} and \S\ref{ch39:hdlt}: we lean
so hard on undirected, spanning-forest-based properties precisely because they are
the ones that slip past the wall. The streaming and sketching world of Chapter~40
attacks the same problems under a different resource (space), and the cut-set
sketch of \S\ref{ch39:worstcase} is the crossover point.
\end{connectionbox}

\section{Putting it together}
\label{ch39:summary}

\begin{complexitybox}
\renewcommand{\arraystretch}{1.25}
\begin{center}
\begin{tabular}{@{}lll@{}}
\toprule
Problem & Model & Per-operation cost \\
\midrule
Connectivity & incremental & $\Oh(\alpha(n))$ amortized (union-find) \\
Connectivity (forest) & fully dynamic & $\Oh(\log n)$ (Euler-tour tree) \\
Connectivity & fully dynamic & $\Oh(\log^2 n)$ amortized (HdLT) \\
Connectivity & fully dynamic, w.c. & $\Oh(\sqrt n)$ det.; polylog rand. (sketch) \\
MST weight & incremental & $\Oh(\log n)$ (link-cut, cycle rule) \\
Min spanning forest & fully dynamic & $\Oh(\log^2 n)$ amortized (leveled ETT) \\
$s$--$t$ reachability (dir.) & fully dynamic & $\Om(n^{1-\eps})$ under OMv \\
\bottomrule
\end{tabular}
\end{center}
\end{complexitybox}

The throughline: maintain a \emph{spanning forest}, store everything else as
spares, and survive deletions by a smaller-half charging argument that bounds how
often any one edge is touched. Euler-tour trees make ``connectivity = same
sequence'' and ``find an incident spare'' both logarithmic; the leveled-forest
wrapper turns the expensive replacement search into a bounded number of
promotions. Where the spanning-forest idea has no purchase -- directed
reachability, exact distances, matching -- the OMv conjecture tells us, with high
confidence, that no polylog method exists, and $\sqrt n$ or $n$ is the best we can
hope for.

\section{Historical Context, Controversies, and the Road Ahead}

\begin{historybox}
In the early decades of algorithm design, graph algorithms were strictly static. If a communication link failed or a new road opened, classic algorithms like Dijkstra's or Tarjan's DFS had to rerun from scratch, consuming linear time per update. The systematic pursuit of dynamic graph algorithms began in $1985$, when Greg Frederickson introduced \emph{topology trees} (hierarchical balanced tree clusterings), achieving $\Oh(\sqrt{m})$ update times for minimum spanning trees. In $1992$, David Eppstein, Zvi Galil, Giuseppe Italiano, and Amnon Nissenzweig developed sparsification trees, reducing the bound to $\Oh(\sqrt{n})$ for general undirected graphs by recursively clustering vertices.

The breakthrough past the square-root barrier arrived in $1995$, when Monika Henzinger and Valerie King introduced randomized Euler-tour trees (ETT, representing trees by their linear traversal sequences in balanced binary search trees) and edge sampling, achieving an amortized update time of $\Oh(\log^3 n)$. In $1998$, Jacob Holm, Kristian de Lichtenberg, and Mikkel Thorup (HdLT) designed the deterministic hierarchical level structure presented here, reaching $\Oh(\log^2 n)$ amortized time per update. The theoretical boundary was solidified in $2015$ when Henzinger, Krinninger, Nanongkai, and Saranurak formulated the Online Matrix-Vector Multiplication ($\mathrm{OMv}$) conjecture (hypothesizing that multiplying an online stream of vectors by a boolean matrix requires essentially cubic work), proving that thirty years of failed attempts to achieve polylogarithmic dynamic reachability on directed graphs reflected a fundamental barrier in fine-grained complexity.

Developing dynamic graph algorithms exposed subtle limitations in standard static data structures. Standard Disjoint-Set Union (DSU) handles insertions in near-constant time, but path compression irrevocably flattens tree structure; undoing operations via a rollback stack works only in offline LIFO order, failing completely when arbitrary edges are deleted online. In addition, when a tree edge is cut, searching blindly through non-tree edges to find a crossing spare can cost $\Omega(m)$ in dense graphs. The HdLT insight was charging replacement searches to edge promotions into the smaller half of the split, guaranteeing that an edge is promoted at most $\log n$ times before reaching its maximum level. What remains an active frontier is achieving deterministic worst-case polylogarithmic update times for general dynamic connectivity against an adaptive adversary.
\end{historybox}

\subsection{What We Can Now Do vs.\ Honest Downsides and the Road Ahead}
\begin{connectionbox}
\textbf{Capabilities Unlocked}:
\begin{itemize}
  \item We can maintain undirected graph connectivity and minimum spanning forests dynamically under arbitrary edge insertions and deletions in $O(\log^2 n)$ amortized time.
  \item Euler-tour trees allow tree operations (\textsc{Link}, \textsc{Cut}, component membership, and subtree aggregates) in $O(\log n)$ worst-case time.
  \item The HdLT hierarchical edge-level framework prevents redundant edge examinations by guaranteeing that every unhelpful edge is permanently promoted to a higher level.
\end{itemize}

\textbf{Honest Downsides and Limitations}:
\begin{itemize}
  \item \emph{Amortized vs.\ worst-case latency}: The $O(\log^2 n)$ bound is strictly amortized. An adversarial deletion can trigger thousands of edge promotions across multiple levels, resulting in occasional $\Omega(n)$ worst-case latency spikes that are unacceptable in hard real-time systems.
  \item \emph{The directed barrier}: Dynamic reachability and shortest paths on directed graphs cannot achieve polylogarithmic update times unless the $\mathrm{OMv}$ conjecture is false; they remain fundamentally constrained to polynomial update times.
  \item \emph{Memory footprint}: Maintaining $O(\log n)$ levels of Euler-tour trees and augmented balanced search trees incurs considerable pointer and memory overhead compared to static adjacency lists.
\end{itemize}

\textbf{Where We Go Next}:
\begin{itemize}
  \item In Chapter~40 (Streaming and Massive Graphs), we consider graphs whose edge sets exceed main memory altogether. When edges arrive as a stream of insertions and deletions, we compress connectivity, cuts, and approximate distances into $\Oh(n\,\poly\log n)$ space via spanners and linear $L_0$-sketches.
  \item In Chapter~41 (Laplacian Solvers and Electrical Flows), we explore how dynamic graph concepts interface with numerical optimization, where dynamic resistance updates accelerate electrical flow computations for maximum flow.
\end{itemize}
\end{connectionbox}

\section{Exercises}
\label{ch39:exercises}

\begin{exercise}[Offline beats online]
Implement offline fully dynamic connectivity using a segment tree over the
timeline plus rollback union-find (union by rank, no path compression, LIFO undo).
Each edge is alive on an interval of operations; place it on the canonical segment
nodes, DFS unioning on entry and rolling back on exit, and answer each query at its
leaf. Prove the total time is $\Oh((m+q)\log q\,\log n)$ and explain precisely why
this technique fails when queries are \emph{online}.
\end{exercise}

\begin{exercise}[Euler-tour subtree sums]
Augment an Euler-tour tree so that each vertex carries a value and you support
\textsc{SubtreeSum}$(v)$ = sum of values in $v$'s subtree, under \textsc{Link},
\textsc{Cut}, and point updates, all in $\Oh(\log n)$. Argue why a subtree is a
contiguous range of the tour, and explain why the same trick does \emph{not} give
you path sums (so you would need a link-cut tree for those).
\end{exercise}

\begin{exercise}[The smaller-half budget]
In HdLT \textsc{Replace}, prove invariant (I2) is preserved: when you promote every
level-$j$ edge inside the smaller side $T$ to level $j{+}1$, every level-$(j{+}1)$
component still has at most $\floor{n/2^{j+1}}$ vertices. Then show that if
you mistakenly scanned and promoted the \emph{larger} half instead, (I2) can fail,
and exhibit a sequence of operations whose amortized cost then degrades to
$\Om(\sqrt n)$.
\end{exercise}

\begin{exercise}[Counting components]
Extend HdLT so that \textsc{NumComponents}() returns the current number of
connected components in $\Oh(1)$. (Hint: a forest on $n$ vertices with $k$ trees
has exactly $n-k$ tree edges; maintain the count of tree edges across all
\textsc{Link}/\textsc{Cut}.) State exactly which lines of
Algorithm~\ref{ch39:alg-hdlt} must update the counter.
\end{exercise}

\begin{exercise}[Incremental MST, the easy side]
Using a link-cut tree with path-max, implement \textsc{InsertMST} (cycle rule) so
that inserting an edge updates the MST in $\Oh(\log n)$ amortized. Prove
correctness from the cycle rule, and explain in one sentence why the \emph{deletion}
side cannot be handled by the link-cut tree alone (what does the cut rule need that
a single tree cannot provide?).
\end{exercise}

\begin{exercise}[Break the weight order]
Construct a small fully dynamic graph and a deletion where the HdLT replacement
search, if it returns the \emph{first} crossing spare it finds rather than the
\emph{minimum-weight} one, produces a spanning forest that is not minimum. Then fix
the search to restore the cut rule and re-verify on your example.
\end{exercise}

\begin{exercise}[The cut-set sketch]
Give each edge a uniformly random $2\ceil{\log n}$-bit id and store at each
vertex the XOR of incident edge ids. Prove that for any vertex set $S$ with exactly
one edge crossing $(S,\bar S)$, the XOR of stored values over $S$ equals that
edge's id. Then design an $\Oh(\log n)$-layer scheme that, with high probability,
recovers a crossing edge even when the number of crossers is unknown. Where does
the failure probability come from?
\end{exercise}

\begin{exercise}[OMv reduction, in full]
Write out the reduction of \S\ref{ch39:omv} for fully dynamic $s$--$t$
reachability completely: specify the gadget graph, the update sequence per query
vector $v_t$, how the $n$ reachability queries reconstruct $Mv_t$, and the cleanup.
Conclude that an $\Oh(n^{1-\eps})$-per-operation structure would refute OMv. Then
identify exactly which step of your reduction breaks if the graph is
\emph{undirected} -- and connect this to why undirected connectivity escapes the
barrier.
\end{exercise}

\begin{exercise}[Decremental only]
Suppose edges are only ever deleted (decremental connectivity), starting from a
known graph. Show how to preprocess so that the \emph{total} time over \emph{all}
deletions plus connectivity queries is $\Oh(m\log^2 n)$, and argue why the
decremental setting can sometimes beat the fully dynamic bound (hint: levels never
need to go back down, so part of the bookkeeping disappears).
\end{exercise}

\begin{exercise}[Worst-case via lazy rebuilding]
Sketch how to convert an amortized $\Oh(\log^2 n)$ structure into one with
$\Oh(\sqrt n)$ worst-case per operation by combining sparsification with periodic
rebuilding of overloaded levels. Identify the trade-off: what do you give up in the
amortized total to buy the worst-case guarantee?
\end{exercise}
